% ===================================================================
%  اللوحة رقم ٦ — البيت من جوّا، الفرش، الأكل، الإنارة.
%
%  Traduction, faite en 2026, en arabe levantin d'aujourd'hui, sur
%  l'ido de texto/io. Regles de la colonne : voir 10-lawha-01.tex.
%  Cinq pieces, cinq numerotations qui repartent a 1 — la chambre, le
%  salon, la salle a manger, la cuisine, la salle de bain. Deux notes
%  d'auteur, t06-noto-1 et t06-noto-2, chacune dans son propre
%  \VUnotes. 62 blocs, 211 renvois : le plus dense du livret.
%
%  LA SALLE DE BAIN (piece 5) EST LE MOT LE PLUS CHARGE DE TOUTE LA
%  COLONNE, ET IL FALLAIT L'ECRIRE SANS RIEN AJOUTER. « الحمّام », en
%  levantin, c'est a la fois la piece d'eau d'un appartement de 1926
%  et le bain public a coupole, celui de Damas et d'Alep, ou l'on va
%  en famille et ou l'on passe l'apres-midi. Le francais a deux mots,
%  « salle de bain » et « hammam », et il a pris le second a cette
%  langue-ci pour ne nommer que le second sens. La colonne ecrit
%  « الحمّام » pour la piece du livret, parce que c'est ce qu'on dit,
%  et le lecteur entend les deux. ON NE PEUT PAS RETIRER A UN MOT CE
%  QU'IL PORTE ; ON PEUT SEULEMENT NE PAS L'AJOUTER SOI-MEME.
%
%  ET LA CUISINE EST L'ENDROIT OU CETTE COLONNE EST CHEZ ELLE. Quatre
%  vingt-dix renvois pour la seule piece 4, et presque aucun a
%  transcrire : الطنجرة la marmite (44), المقلاية la poele (47),
%  الغربال le tamis (28), المبشرة la rape (29), المغرفة la louche
%  (21), الهاون, البلاّعة, الكاز. La colonne gujaratie avait eu son
%  tableau chez elle au port ; celle-ci l'a dans une cuisine, et pour
%  la meme raison — c'est ce qu'on manie tous les jours qui garde son
%  nom.
%
%  TROIS MOTS DE PLUS POUR LE COMPTE DU MANDAT, et ils sont tous du
%  meme cote de la piece : l'abat-jour (13) est « الأباجورة » comme au
%  tableau 1, le lustre (38) est « النجفة » — qui est arabe, celui-la,
%  et vient de Najaf — et le petrole (12) est « الكاز », de l'anglais
%  gas par le turc. CE QUI ECLAIRE EST TOUJOURS EMPRUNTE, DANS CETTE
%  COLONNE, ET CE QUI CHAUFFE AUSSI : la صوبيا du tableau 1 etait
%  turque. Ce qui se mange et ce qui se lave sont restes arabes.
%
%  Enfin le paravent (36) et le vase de Chine (34) viennent de Chine
%  dans le texte meme, et leurs noms viennent d'ailleurs encore : le
%  premier est « البارافان », francais, le second « الفازة » — et la
%  porcelaine est « الصيني », la chinoise. LE LIVRET FAIT VENIR UN
%  OBJET DE CHINE PAR UN ONCLE ; LA LANGUE, ELLE, LE FAIT VENIR PAR
%  TROIS PAYS.
% ===================================================================

%%K t06-apar-1 apar
\VUsaut{6.0mm}
\VUtitre{15.0pt}{60}{اللوحة رقم ٦}
\VUsaut{1.4mm}
\VUtitre{11.4pt}{0}{البيت من جوّا، الفرش، الأكل،}
\VUsaut{0.4mm}
\VUtitre{11.4pt}{0}{الإنارة.}
\VUsaut{2.2mm}

%%K t06-apar-2 apar
البيت خلص. وعمّ حنّا انتقل عبيته الجديد، اللي منعرف ترتيبه. هلّق
خلّينا نشوف كيف فرشه. رح نتفرّج أول شي على :

%%K t06-tit-1 sub
\VUpk{10.2pt}{40}{غرفة النوم.}
\VUsaut{3.0mm}

%%K t06-c1-01-1 p
1. --- وصلنا بوقت مش مناسب، لأن الصبيّة عم تنضّف \VUgras{الغرفة}.

%%K t06-c1-02-1 p
2. --- على \VUgras{طاولة الزينة} \textsuperscript{(1)} في هون وهونيك أغراض
مختلفة بتنغسل فيها ويتمشّط فيها الواحد، يعني بيعمل زينته. وهنّي
\VUgras{الطشت} \textsuperscript{(2)} و\VUgras{إبريق الميّ} \textsuperscript{(3)}،
و\VUgras{فرشاية الراس} \textsuperscript{(4)}، و\VUgras{المشط}
\textsuperscript{(5)}، و\VUgras{الصابون} \textsuperscript{(6)}
و\VUgras{الإسفنجة} \textsuperscript{(7)}، وأخيراً \VUgras{قنّينة زغيرة}
\textsuperscript{(8)} لمعجون السنان السايل.

%%K t06-c1-03-1 p
3. --- وعالأرض، قدّام طاولة الزينة، هيّو \VUgras{السطل}
\textsuperscript{(9)} لجمع الميّ الوسخة.

%%K t06-c1-04-1 p
4. --- ومنشوف بـ\VUgras{غرفة المدخل} \textsuperscript{(10)} كم
\VUgras{علّاقة} \textsuperscript{(13)} وشمسيّتين \textsuperscript{(11)}
بـ\VUgras{حاملة الشماسي} \textsuperscript{(12)}.

%%K t06-c1-05-1 p
5. --- وعالجهة التانية من الباب في \VUgras{خزانة المراية}
\textsuperscript{(14)} اللي فيها \VUgras{البياضات} \textsuperscript{(15)}.

%%K t06-c1-06-1 p
6. --- خلّينا نستفيد وقت \VUgras{الصبيّة} \textsuperscript{(6)} عم ترتّب
\VUgras{التخت} \textsuperscript{(17)}، لنتفرّج على أقسامه. \VUgras{هيكل
التخت} \textsuperscript{(20)} فيه \VUgras{فرشة زنبركات}
\textsuperscript{(21)} منيحة، وعليها يوستينا حالاً رح تحطّ \VUgras{فرشة}
\textsuperscript{(22)} صوف. وبعدين رح تاخد من الكراسي \VUgras{الشرشفين}
\textsuperscript{(23)} و\VUgras{الحرام} \textsuperscript{(24)}. وبعدها رح تحطّ
عند راس التخت \VUgras{المخدّة العريضة} \textsuperscript{(25)}
و\VUgras{المخدّة} \textsuperscript{(26)} اللي مع \VUgras{اللحاف}
\textsuperscript{(27)} عالـ\VUgras{سجّادة} \textsuperscript{(32)}. وعم تستعمل
منفضة ريش لتشيل الغبرة عن \VUgras{البردايات} \textsuperscript{(19)}
و\VUgras{ناموسية التخت} \textsuperscript{(18)}، وهيك بيكون قسم من الشغل
خلص.

%%K t06-c1-07-1 p
7. --- وبعدين لازم ترتّب شوي \VUgras{كومودينو التخت}
\textsuperscript{(28)}، وتعبّي \VUgras{المنبّه} \textsuperscript{(29)}،
وتحضّر \VUgras{لمبة الليل} \textsuperscript{(30)}، وتشيل \VUgras{الشمعة}
\textsuperscript{(31)} لتنضّف الشمعدان.

%%K t06-tit-2 sub
\VUsaut{5.0mm}
\VUpk{10.2pt}{40}{سلّة الشغل وعدّة المدفأة.}
\VUsaut{3.0mm}

%%K t06-c1-08-1 p
8. --- و\VUgras{سلّة الشغل} \textsuperscript{(34)} اللي جنب
\VUgras{الكرسي الزغير} \textsuperscript{(33)} كمان كتير بدّها ترتيب. لازم
تطوي \VUgras{الأشغال المطرّزة} \textsuperscript{(35)}، وتحطّ
بـ\VUgras{طاولة الشغل} \textsuperscript{(38)} \VUgras{الكشتبان}
و\VUgras{الخيطان} \textsuperscript{(36)} و\VUgras{الإبر} \textsuperscript{(37)}
و\VUgras{المقصّ} \textsuperscript{(39)}، وتسكّر بـ\VUgras{المفتاح}
\textsuperscript{(40)} غطا الطاولة.

%%K t06-c1-09-1 p
9. --- وعالـ\VUgras{طاولة أبو رجل وحدة} \textsuperscript{(41)} اللي
بالنص، يوستينا رح تغيّر ميّ \VUgras{الكاراف} \textsuperscript{(42)}. ورح
تحطّ \VUgras{سكّر} بـ\VUgras{السكّرية} \textsuperscript{(43)}، وتنضّف
\VUgras{ملقط السكّر} \textsuperscript{(44)}، وتشيل \VUgras{فرشاية التياب}
\textsuperscript{(46)}.

%%K t06-c1-10-1 p
10. --- والجهة اليمنى من الغرفة رح تطلب منها شغل أقل، لأن
\VUgras{الكنباية الطويلة} \textsuperscript{(47)} و\VUgras{مخدّتها}
\textsuperscript{(48)} بمحلّهن الصحّ، ولأنه الدني مش باردة، ما في شي
مزحزح بين عدّة \VUgras{المدفأة} \textsuperscript{(49)}.
\VUgras{المجرفة} \textsuperscript{(50)} و\VUgras{الملقط} \textsuperscript{(51)}
و\VUgras{الحمّالة} \textsuperscript{(55)} و\VUgras{حاجز الرماد}
\textsuperscript{(56)} كلهن نضاف ؛ و\VUgras{حاجز النار}
\textsuperscript{(54)} ما انزحزح، و\VUgras{صندوق الحطب}
\textsuperscript{(52)} مليان \VUgras{حطب} \textsuperscript{(53)}.

%%K t06-noto-1 noto
\VUnotes{143.2mm}{%
(*) Babo = طفل زغير، بالإنكليزي \textit{baby}، وبالفرنسي
\textit{bébé}، وبالإيطالي \textit{bambino}.%
}

%%K t06-noto-2 noto
\VUnotes{143.2mm}{%
(*) Lambrequino = بالفرنسي \textit{lambrequin}، وبالإسباني
\textit{lambrequines}، وبالبرتغالي \textit{lambrequins}، وحتى
بالألماني والإنكليزي \textit{lambrequins} --- يعني وحدة بالألماني
والإنكليزي والفرنسي والبرتغالي والإسباني.%
}

%%K t06-c1-11-1 p
11. --- ومع هيك لازم تشيل الغبرة عن \VUgras{ساعة الحيط}
\textsuperscript{(57)}، و\VUgras{الشمعدانات} \textsuperscript{(58)}،
و\VUgras{علب الحاجيات} \textsuperscript{(59)}، وأخيراً تمسح
\VUgras{المراية} \textsuperscript{(60)}.

%%K t06-c1-12-1 p
12. --- أما \VUgras{مهد} \textsuperscript{(61)} الولد الزغير (الـ
babo (*))، فهوّي جاهز من قبل.

%%K t06-tit-3 sub
\VUsaut{5.0mm}
\VUpk{10.2pt}{40}{الصالون.}
\VUsaut{3.0mm}

%%K t06-c2-01-1 p
1. --- وهلّق هيّو الصالون. \VUgras{غرفة} كتير حلوة. والمدفأة، اللي
بالزاوية اليمنى، مزيّنة بـ\VUgras{تمثال زغير} \textsuperscript{(1)} ناعم،
وبـ\VUgras{فازتين} \textsuperscript{(3)} حلوين، ومكسيّة بـ\VUgras{ستارة
مخمل} \textsuperscript{(2)} (*).

%%K t06-c2-02-1 p
2. --- ولمنع شرار الموقدة يحرق السجّادة، حطّوا قدّام الموقدة
\VUgras{حاجز شرار} \textsuperscript{(4)} من نحاس.

%%K t06-c2-03-1 p
3. --- وجنبه، \VUgras{مكتب} \textsuperscript{(5)} نسائي أنيق مفتوح. عليه
\VUgras{ستّ البيت} \textsuperscript{(23)} عم تكتب مكاتيبها.

%%K t06-c2-04-1 p
4. --- والشبّاك مزيّن بـ\VUgras{بردايات إزاز} \textsuperscript{(6)} مع
\VUgras{رباطات} \textsuperscript{(8)} معلّقة بـ\VUgras{المشاجب}
\textsuperscript{(7)}، وببردايات قماش \VUgras{مسدّلة من فوق}
\textsuperscript{(9)}.

%%K t06-c2-05-1 p
5. --- وبحضن الشبّاك، \VUgras{حاملة قصريّة} \textsuperscript{(11)} فيها
\VUgras{نبتة} \textsuperscript{(10)} خضرا، نخلة حلوة ورقها عم يمتدّ لعند
\VUgras{لمبة الكاز} \textsuperscript{(12)} تقريباً.

%%K t06-c2-06-1 p
6. --- و\VUgras{الأباجورة} \textsuperscript{(13)} تبع اللمبة رتّبتها
ماريا، الصبيّة \textsuperscript{(14)} الحلوة اللي عند البيانو
\textsuperscript{(15)}. قاعدة كتير معتدلة على \VUgras{الكرسي الزغير}
\textsuperscript{(16)}، عم تدرس قطعة موسيقى. وهيّي كتير شاطرة ومرتّبة،
ودليل هيك \VUgras{خزانة الموسيقى} \textsuperscript{(17)} تبعها اللي مرتّبة
تماماً.

%%K t06-tit-4 sub
\VUsaut{5.0mm}
\VUpk{10.2pt}{40}{الشقي.}
\VUsaut{3.0mm}

%%K t06-c2-07-1 p
7. --- وأخوها، باولوس، عم يدوّر على \VUgras{كتاب} \textsuperscript{(19)}
بـ\VUgras{المكتبة} \textsuperscript{(18)}. هوّي \VUgras{شاب}
\textsuperscript{(20)}، ما بيقعد بمحلّه وما بينحمل. وهوّي عم يتعلّق
بالمكتبة ليوصل للكتاب اللي كتير فوق، عم يخاطر يوقّع
\VUgras{التمثال النصفي} \textsuperscript{(21)} اللي فوق.

%%K t06-c2-08-1 p
8. --- وعم يستغلّ، ليعمل هالحماقة، إنّه إمّه مشغولة عم تحكي مع صاحبتها،
\VUgras{سيّدة} \textsuperscript{(22)} إجت تقضّي كم يوم بالبيت. والإم قاعدة
على \VUgras{الكنباية} \textsuperscript{(24)} جنب \VUgras{الفوتيه}
\textsuperscript{(25)}، وصاحبتها على \VUgras{الكرسي} \textsuperscript{(27)}
اللي ما منشوف منه غير \VUgras{الظهر} \textsuperscript{(26)}.

%%K t06-c2-09-1 p
9. --- و\VUgras{البرواز} \textsuperscript{(30)} الكبير المعلّق عالحيط فيه
\VUgras{صورة} \textsuperscript{(29)} قريبة. وبـ\VUgras{الألبوم}
\textsuperscript{(32)} المحطوط عالطاولة مجموعة \VUgras{صور}
\textsuperscript{(33)} باقي أهل البيت. والولاد هلّق كانوا عم يقلّبوه،
لأن \VUgras{الكراسي} \textsuperscript{(31)} لسّاتها مجمّعة حوالين الطاولة.

%%K t06-c2-10-1 p
10. --- و\VUgras{الفازة الصينية} \textsuperscript{(34)} من خزف، اللي جنب
\VUgras{سكّينة الورق} \textsuperscript{(35)}، انهدت لماريا بعيدها،
و\VUgras{البارافان} \textsuperscript{(36)} اللي مزيّن زاوية الغرفة جابه
عمّها من الصين.

%%K t06-c2-11-1 p
11. --- ومع \VUgras{اللوحات} \textsuperscript{(37)} الحلوة المعلّقة
عالـ\VUgras{حيط} \textsuperscript{(42)}، ومع ضو \VUgras{النجفة}
\textsuperscript{(38)} اللي \VUgras{لمباتها الكهربائية}
\textsuperscript{(39)} عم تلمع بفرح بالليل، هالصالون بيبيّن كتير مريح.
و\VUgras{السجّادة} \textsuperscript{(40)} الطريّة المفروشة عالباركيه،
و\VUgras{الجرس الكهربائي} \textsuperscript{(41)} السهل الاستعمال، بيكمّلوا
راحته.

%%K t06-tit-5 sub
\VUsaut{3.6mm}
\VUpk{10.2pt}{40}{غرفة السفرة.}
\VUsaut{3.0mm}

%%K t06-c3-01-1 p
1. --- دقّ جرس العشا. وكل العيلة، ما عدا ماريا، مجتمعة حوالين
الطاولة. وغرفة السفرة دافية منيح بفضل \VUgras{الصوبيا}
\textsuperscript{(1)} الممتازة من خزف، مع \VUgras{بوري}
\textsuperscript{(2)} رفيع.

%%K t06-c3-02-1 p
2. --- وهالغرفة ما فيها فرش كتير. عالحيط معلّقة، فوق \VUgras{الرفّ}
\textsuperscript{(4)}، \VUgras{فسقيّة} \textsuperscript{(3)} للزينة. وكتير
قريب، في \VUgras{البوفيه} \textsuperscript{(5)} اللي بيحطّوا عليه الأكل
البارد، وصحون الحلو، وأدوات السفرة اللي ما بدّهن ياها حالاً.

%%K t06-c3-03-1 p
3. --- وهيك منشوف عالرفّ الفوقاني \VUgras{فرّوج مشوي}
\textsuperscript{(7)}، و\VUgras{سمك} \textsuperscript{(8)}، و\VUgras{طاسة}
\textsuperscript{(10)} فيها \VUgras{فواكه} \textsuperscript{(9)}. وتحت هيّو
\VUgras{الجبنة} \textsuperscript{(11)}، و\VUgras{الخبز} \textsuperscript{(12)}،
و\VUgras{حاملة القناني} \textsuperscript{(13)} اللي فيها كل شي بدّه ياه
السلطة : الزيت، والخلّ، و\VUgras{الملح} \textsuperscript{(14)}
و\VUgras{البهار} \textsuperscript{(15)}. وأخيراً، عالرفّ التحتاني في
\VUgras{كاتو} \textsuperscript{(17)} جنب \VUgras{وعا الخردل}
\textsuperscript{(16)}.

%%K t06-c3-04-1 p
4. --- وساعة غرفة السفرة، اللي بيسمّوها \VUgras{ساعة الحيط}
\textsuperscript{(20)}، إلها شكل كتير خاص. مركّبة بإطار خشب فيه كمان
\VUgras{بارومتر} \textsuperscript{(19)} و\VUgras{ترمومتر}
\textsuperscript{(18)}.

%%K t06-tit-6 sub
\VUsaut{4.6mm}
\VUpk{10.2pt}{40}{خدمة العشا.}
\VUsaut{3.0mm}

%%K t06-c3-05-1 p
5. --- وعكل حيطان الغرفة مركّب \VUgras{لبيس خشب} \textsuperscript{(21)}
من سنديان مشمّع. و\VUgras{ستارة الباب} \textsuperscript{(22)} من قماش
بتسكّر منيح الباب اللي بيفوت عالفرندة. وهونيك رح تروح العيلة تشرب
القهوة بعد العشا. ومن هلّق \VUgras{الفناجين} \textsuperscript{(24)}
و\VUgras{الصحون الزغيرة} \textsuperscript{(25)} جاهزين على
\VUgras{النملية} \textsuperscript{(23)}.

%%K t06-c3-06-1 p
6. --- وفوق الطاولة مثبّت بالسقف \VUgras{ثريّا نازلة}
\textsuperscript{(26)} لمبتها كتير قوية وبتنوّر بالليل كل غرفة السفرة.

%%K t06-c3-07-1 p
7. --- وهلّق خلّينا نتفرّج على \VUgras{طاولة الأكل} \textsuperscript{(27)}
نفسها. لمّا فاتت العيلة، كانت السفرة جاهزة. وعالـ\VUgras{شرشف}
\textsuperscript{(28)} كان محطوط، بمحلّ كل واحد، \VUgras{صحن}
\textsuperscript{(33)}، و\VUgras{فوطة} \textsuperscript{(29)}،
و\VUgras{ملعقة} \textsuperscript{(34)}، و\VUgras{شوكة} \textsuperscript{(35)}،
و\VUgras{سكّينة} \textsuperscript{(36)}، و\VUgras{كاسة}
\textsuperscript{(32)}. وكان في كمان \VUgras{قناني} \textsuperscript{(30)}
للنبيد و\VUgras{كاراف} \textsuperscript{(31)} للميّ.

%%K t06-c3-08-1 p
8. --- و\VUgras{الخادم} \textsuperscript{(39)} أول شي جاب \VUgras{وعا
الشوربة} \textsuperscript{(37)}، اللي لسّا فيه شوربة عم تبخّر. وهلّق عم
يقدّم \VUgras{الصحن} \textsuperscript{(38)} الأول، صحن اللحمة. وبعدين رح
يقدّم اللي سمّيناهن، وأخيراً، بعد \VUgras{صحن الحلو}، رح يستفيد إنّه
أصحاب البيت راحوا عالفرندة، ليشيل أدوات السفرة ويرتّب غرفة السفرة.

%%K t06-c3-08-2 p
\textit{ملاحظة.} --- \textit{الكلمات الشايعة اللي إلها علاقة بالأكل
مذكورة هون بشكل مختصر كتير. والقائمة رح تنكمّل باللوحتين « الفندق »
(رقم ١٣) و« السوق » (رقم ١٥).}

%%K t06-tit-7 sub
\VUsaut{4.6mm}
\VUpk{10.2pt}{40}{المطبخ.}
\VUsaut{3.0mm}

%%K t06-c4-01-1 p
1. --- وبالمطبخ، اللي عم نتطلّع عليه من الباب النصّ مفتوح، منشوف فوضى
حلوة. قدّام \VUgras{الموقدة} \textsuperscript{(1)} اللي عم يفرقع فيها
\VUgras{النار} \textsuperscript{(3)}، مركّب \VUgras{دوّار السيخ}
\textsuperscript{(5)}. و\VUgras{مشوي} \textsuperscript{(7)} حلو
\VUgras{مسيّخ} \textsuperscript{(6)} وعم ينشوي.

%%K t06-c4-02-1 p
2. --- و\VUgras{المنفاخ} \textsuperscript{(8)} و\VUgras{مجرفة الفحم}
\textsuperscript{(9)}، اللي هلّق استعملوهن، ضلّوا عالأرض متل
\VUgras{الحطب} \textsuperscript{(10)}.

%%K t06-c4-03-1 p
3. --- وفوق المدفأة مرتّب ؛ \VUgras{الفانوس} \textsuperscript{(11)}،
و\VUgras{علبة الكبريت} \textsuperscript{(13)} اللي فيها
\VUgras{الكبريت} \textsuperscript{(12)}، و\VUgras{الشمعدان}
\textsuperscript{(14)} و\VUgras{حاملة الشمعة} \textsuperscript{(15)} مع
\VUgras{الشمعة} \textsuperscript{(16)}، كلهن بمحلّهن. بس \VUgras{سطل
الفحم} \textsuperscript{(18)} الكبير، المليان \VUgras{فحم حجري}
\textsuperscript{(17)} اللي \VUgras{الطبّاخة} \textsuperscript{(23)} هلّق
عبّته من القبو، ضلّ بنص المطبخ.

%%K t06-c4-04-1 p
4. --- والحقيقة، هالمرا المسكينة لازم تستعجل ليكون الغدا جاهز بوقته.
هلّق هيّي جنب \VUgras{طبّاخة الغاز} \textsuperscript{(19)} وعم تحرّك
\VUgras{صلصة} \textsuperscript{(24)} لازم ما تنحرق كتير.

%%K t06-c4-05-1 p
5. --- وبـ\VUgras{الفرن} \textsuperscript{(20)} عم ينخبز كاتو حضّرته
لتصبيرة الولاد. وفوق طبّاخة الغاز معلّقين \VUgras{المغرفة}
\textsuperscript{(21)} و\VUgras{مغرفة الزبد} \textsuperscript{(22)}.

%%K t06-tit-8 sub
\VUsaut{4.0mm}
\VUpk{10.2pt}{40}{العدّة.}
\VUsaut{3.0mm}

%%K t06-c4-06-1 p
6. --- و\VUgras{عدّة الطبخ} \textsuperscript{(25)} المعلّقة بآخر الغرفة
كتير نضيفة. \VUgras{الطناجر} \textsuperscript{(26)} النحاس الحلوة،
و\VUgras{وعا الحليب} \textsuperscript{(27)}، و\VUgras{الغربال}
\textsuperscript{(28)} و\VUgras{المبشرة} \textsuperscript{(29)} كلهن
انضّفوا بعناية.

%%K t06-c4-07-1 p
7. --- و\VUgras{علبة الملح} \textsuperscript{(30)} محطوطة عالنملية جنب
\VUgras{إبريق الشاي} \textsuperscript{(31)}، و\VUgras{قنّينة الزيت}
\textsuperscript{(32)} الكبيرة. و\VUgras{الجارور} \textsuperscript{(33)}
يمكن فيه أدوات السفرة وسكاكين المطبخ.

%%K t06-c4-08-1 p
8. --- و\VUgras{المكنسة} \textsuperscript{(34)} و\VUgras{منفضة الريش}
\textsuperscript{(35)} بالزاوية. والطبّاخة هلّق استعملت \VUgras{لوح
الخشب} \textsuperscript{(36)}، وتركته جنب الشبّاك.

%%K t06-c4-09-1 p
9. --- و\VUgras{منشفة الإيدين} \textsuperscript{(37)} معلّقة عالحيط، جنب
\VUgras{القمع} \textsuperscript{(38)}. وبهالقسم من المطبخ متروكين
\VUgras{الشوّاية} \textsuperscript{(39)}، و\VUgras{الساطور}
\textsuperscript{(40)} و\VUgras{لوح الفرم} \textsuperscript{(41)}. وبعدين،
فوق الساطور، على \VUgras{رفّ} \textsuperscript{(42)}، في \VUgras{أوعية}
\textsuperscript{(43)}، و\VUgras{وعا السلق} \textsuperscript{(44)} مع
\VUgras{غطاه} \textsuperscript{(45)} و\VUgras{الطنجرة}
\textsuperscript{(46)}. و\VUgras{المقلاية} \textsuperscript{(47)} معلّقة
جنبهن.

%%K t06-c4-10-1 p
10. --- وبس \VUgras{الحنفية} \textsuperscript{(48)} النحاس عم تلمع
بهالزاوية. و\VUgras{المجلى} \textsuperscript{(49)} اللي محطوط عليه
\VUgras{الجرّة} \textsuperscript{(50)} التخينة، يمكن كتير مريح مع خزانته
الزغيرة، اللي بينحفظ فيها \VUgras{برميل الخلّ} \textsuperscript{(51)} مع
\VUgras{حنفيته} \textsuperscript{(52)}، و\VUgras{علبة الزبالة}
\textsuperscript{(53)} و\VUgras{الصحون الوسخة} \textsuperscript{(54)}.

%%K t06-tit-9 sub
\VUsaut{4.0mm}
\VUpk{10.2pt}{40}{أغراض تانية محطوطة بالمطبخ.}
\VUsaut{3.0mm}

%%K t06-c4-11-1 p
11. --- و\VUgras{اللقّان} \textsuperscript{(55)} اللي بينغسل فيه
\VUgras{الخضرة} \textsuperscript{(58)}، و\VUgras{الدست}
\textsuperscript{(56)} المصبوب المحطوط عالـ\VUgras{بلاط}
\textsuperscript{(57)}، يمكن كمان إلهن محلّ بهالخزانة.

%%K t06-c4-12-1 p
12. --- والطبّاخة أكيد ما ناقصها مواقد، لأنه هيّو كمان، على طرف
الطاولة، \VUgras{موقد غاز} \textsuperscript{(59)} عم تسخن عليه ميّ
بطنجرة زغيرة. و\VUgras{البخار} \textsuperscript{(60)} اللي عم يطلع منها
بيدلّنا إنّها يمكن عم تغلي. والأرجح إنّه هالميّ رح تنستعمل للقهوة،
لأنه هيّن عالطرف التاني من الطاولة، \VUgras{ركوة القهوة}
\textsuperscript{(64)} و\VUgras{مطحنة القهوة} \textsuperscript{(63)}.

%%K t06-c4-13-1 p
13. --- والموقد موصول بـ\VUgras{عين الغاز} \textsuperscript{(62)}
بـ\VUgras{خرطوم مطّاط} \textsuperscript{(61)} طويل.

%%K t06-c4-14-1 p
14. --- و\VUgras{الخدّامة اليومية} \textsuperscript{(66)} اللي بتيجي
تساعد الطبّاخة، عم تحضّر الخضرة للعشا. ولمّا تخلص تنظّفها وغسيلها، رح
تحطّها بـ\VUgras{الطاسة} \textsuperscript{(65)} لحدّ ما يجي وقت طبخها.

%%K t06-tit-10 sub
\VUsaut{4.0mm}
{\centering\fontsize{10.2pt}{10.2pt}\selectfont\textls[40]{\scshape الحمّام.} \textit{(غرفة الاستحمام.)}\par}
\VUsaut{3.0mm}

%%K t06-c5-01-1 p
1. --- ما ضلّ إلنا غير قلّة أدب وحدة لنعرف أهم غرف الشقّة : نشوف
تركيبة الحمّام. وهاد ما رح ياخد وقت طويل، لأنه مش كبير، وبسرعة رح
نعرف إنّه \VUgras{بانيو} \textsuperscript{(1)} فيه \VUgras{مسخّن}
\textsuperscript{(2)} منيح، وإنّه \VUgras{صحن الصابون} \textsuperscript{(3)}
بيوصلّه اللي عم يستحمّي بإيده، وإنّه \VUgras{حاملة المناشف}
\textsuperscript{(5)} أكيد بتخلّي \VUgras{مناشف الزينة}
\textsuperscript{(4)} تنشف بسهولة.

%%K t06-c5-02-1 p
2. --- وهالغرفة حيطانها مكسيّة بـ\VUgras{القيشاني} \textsuperscript{(6)}،
اللي خلّاهن يقدروا يركّبوا \VUgras{دوش} \textsuperscript{(7)} بلا ما
يخافوا إنّه الميّ، اللي بترشّ لورا وقت استعماله، تخرّب الحيطان.
و\VUgras{طشت} \textsuperscript{(8)} زنك بيخدم لحمّام الإجرين، بيكمّل
الفرش.

\VUsaut{6.0mm}
\VUfilet{22mm}
